\subsection{Fixed-rank representations with tame original boundary}
\begin{proposition}[Proved: Boundedness of the canonical parabolic lattices]
\label{D:tame-parabolic-boundedness}
Let $S$ be a normal proper connected surface over an algebraically closed field $k$ of characteristic $p>0$, and let $D$ be a reduced divisor. Fix a finite set $\Sigma$ containing the singular points of $S$, the singular and intersection points of $D$, and any additional omitted points. Put $S^\circ=S\setminus\Sigma$, $U^\circ=S^\circ\setminus D$, and let $Q^\circ$ be the quotient of $\pi_1(U^\circ)$ classifying covers tame at the generic points of the original components of $D$. Fix a projective resolution $\pi:Y\to S$, an isomorphism over $S^\circ$, and a positive exceptional divisor $Z$ with $Z\cdot E_i<0$ for every exceptional component.
For every fixed rank $r$, one bounded family of vector bundles on $Y$ contains minimal extensions of all canonical reflexive lattices of representations $Q^\circ\to\mathrm{GL}_r(\mathbf F_{p^f})$, all their coefficient-Frobenius conjugates, and all distinct parabolic weight lattices between $E(-D_i)$ and $E$, for every $f\ge1$ and every component $D_i$. The family is independent of the Kummer denominators.
\end{proposition}
\begin{proof}
Write $q=p^f$. For a representation $M$, let $M_a$ be its coefficient-Frobenius conjugates, indexed by $a\in\mathbf Z/f\mathbf Z$, and let $E_a$ be their canonical reflexive lattices on $S$. Thus $E_a|_{U^\circ}=\mathcal O_{U^\circ}\otimes_{\mathbf F_q}M_a$. At a generic boundary DVR we choose a common prime-to-$p$ Kummer denominator $n$. The integral invariant lattice has weight representatives $0\le j<n$, and its Frobenius map has exponent $c_j=\lfloor pj/n\rfloor\le p-1$. We obtain injections
\[\phi_a:F_S^{[*]}E_a\longrightarrow E_{a+1}\]
with pure one-dimensional cokernels $T_a$ supported on $D$.
The canonical-lattice estimates proved above remain uniform for this finite cycle. Indeed, their local ordinary/reflexive Frobenius defect ideal and relative-Frobenius determinant depend only on the fixed surface charts. They apply equally to a map from the $a$th Kummer module to the $(a+1)$st. The normalized eigenpiece maps now run around finite cycles of pairs $(a,j)$. Their one-step upper-slope inequalities have the same fixed error, so taking the maximum around such a cycle gives the same upper-slope bound. The fixed matrix comparison and the additive normalization-index bound for the final cokernel grades are unchanged. Consequently
\[\chi(S,T_a)\le C r\]
for one constant independent of $r,f,n,a$. The same argument gives an upper Euler bound $C\operatorname{rk}G$ for each pure whole-eigenpiece grade $G$, with the sum bounded by $Cr$ along each boundary component.
Taking determinants gives
\[c_1(E_{a+1})-p c_1(E_a)=[T_a]_1,\]
where every coefficient of $[T_a]_1$ lies between $0$ and $(p-1)r$. Iterating the $f$-cycle shows
\[c_1(E_a)=-\sum_i\theta_{a,i}[D_i],\qquad 0\le\theta_{a,i}\le r\]
in the numerical divisor space. In particular $K_S\cdot c_1(E_a)=O(r)$. The uniform normal-surface Riemann--Roch theorem \citep[Theorem~0.1]{C_Langer2026} and Frobenius compatibility of the second Chern character \citep[Theorem~4.5(2)]{C_Langer2025} give
\[\operatorname{ch}_2(E_{a+1})-p^2\operatorname{ch}_2(E_a)\le Cr.\]
After cycling, division by $p^{2f}-1$ gives $\operatorname{ch}_2(E_a)\ge-Cr$, uniformly in $f$.
For clarity, the fixed-resolution sandwich has the adjacent-index form
\[F_{a+1}(-A)\subset F_Y^*F_a\subset F_{a+1},\qquad F_a=\pi^{[*]}E_a,\]
for one fixed effective divisor $A$. It follows from the same fixed boundary ideal and stable-MCM Frobenius-defect ideal used in the canonical-lattice proof. Taking the maximum around the finite cycle, both on $Y$ and on a fixed smooth very ample curve not contained in the determinant-zero divisors, gives the same bounds $h^0(S,E_a)\le r$ and $h^2(S,E_a)\le Cr$. Riemann--Roch therefore yields $|\chi(S,E_a)|\le Cr$.
For a weight lattice $L$ with $E_a(-D_i)\subset L\subset E_a$, the successive pure grades of $E_a/L$ have total upper Euler bound $Cr$. Hence $\chi(S,L)\ge-Cr$. We also have $h^0(S,L)\le r$ and
\[\operatorname{Hom}(L,\omega_S)\hookrightarrow\operatorname{Hom}(E_a(-D_i),\omega_S).\]
Principalize the fixed coherent ideal $\mathcal O_S(-D_i)$ on the resolution. The latter Hom space injects into sections of $F_a^\vee$ tensored by one fixed line bundle, so the preceding dual-slope argument bounds it by $Cr$. Thus $|\chi(S,L)|\le Cr$ as well. Its first Chern class differs from that of $E_a$ by a multiple of $D_i$ between $0$ and $r$.

Choose minimal extensions of these reflexive lattices on $Y$. Their restrictions to the exceptional components lie in the fixed slope windows established earlier \citep[Definition~1.2 and Lemma~1.3]{C_Langer2026}. The uniform relative Euler-characteristic estimate $\chi(\pi,F)\le Cr$ therefore bounds $|\chi(Y,F)|$ by $Cr$. The numerical first Chern classes on $S$ lie in a fixed $O(r)$ divisor span; the exceptional slope windows and the negative-definite intersection matrix then bound $c_1(F)\cdot H$ by $Cr$, for a fixed very ample $H$ on $Y$.
The bounded-pole lemma extends each one-step Frobenius map between minimal extensions to
\[F_Y^*\widetilde E_a\hookrightarrow\widetilde E_{a+1}(BZ)\]
with one fixed $B$. Here the restriction-slope bounds for a Frobenius pullback follow from the already used curve Frobenius estimate; all constants are independent of the rank and of $f$. Cycling gives a uniform upper bound on $\mu_{\max,H}(\widetilde E_a)$. The inclusion $L\hookrightarrow E_a$ likewise extends between their minimal extensions with one fixed exceptional pole, and bounds the maximal slope of the extension of $L$. At fixed rank, the resulting Hilbert polynomials range over a finite set. Langer's boundedness theorem \citep[Theorem~4.2, p.~269]{C_LangerBounded2004}, retained by its addendum \citep{C_LangerAddendum2004}, now gives the asserted bounded family.
\end{proof}

\begin{theorem}[Proved: Fixed-rank growth for the original-boundary tame quotient]
\label{D:tame-parabolic-fixed-rank-growth}
In the preceding situation, for every fixed positive integer $r$ there is a constant $C_r$ such that
\[\#\bigl(\operatorname{Hom}_{\mathrm{cont}}(Q^\circ,\mathrm{GL}_r(\mathbf F_{p^f}))/\mathrm{GL}_r(\mathbf F_{p^f})\bigr)\le p^{C_r f}\qquad(f\ge1).\]
The same bound holds for the original quotient before enlarging the finite omitted set. No uniform quadratic dependence on $r$ is asserted.
\end{theorem}
\begin{proof}
We first construct a finite-type space of unweighted parabolic states. A state consists of a reflexive lattice $E$ and, for each original boundary component $D_i$, an ordered chain
\[E=E_{i,0}\supset E_{i,1}\supset\cdots\supset E_{i,l_i}=E(-D_i),\qquad l_i\le r,\]
whose successive ranks on $D_i$ are positive. On $S^\circ$, this is a flag of subbundles of $E|_{D_i^\circ}$; we retain the full chain, including its extension data, and do not yet assign numerical weights.
The preceding proposition supplies bounded minimal extensions for every lattice occurring in a representation. We take finite unions of their parameter spaces and parameterize the chain maps by bounded-pole Hom schemes on the projective surface $Y$. Conditions on support, endpoint, inclusions, and local freeness of the successive quotients on $S^\circ$ are constructible: they are finite-presentation conditions, after the usual finite flattening stratifications, and their failure loci project to constructible sets. An exact isomorphism of states is an isomorphism of $E$ preserving every member of every chain. The fixed-pole lemma and the inverse-product equations used earlier make this an effective constructible equivalence relation on a finite-type parameter set. Beke's effective quotient theorem \citep[Theorem~2; proof via Proposition~3.1]{Beke2011FixedRankB} therefore gives a constructible quotient $P_r$ whose points are the state-isomorphism classes. We may enlarge the parameter family; only the class set of cycles will matter below.
There are finitely many grade-rank types. A one-step Frobenius type consists, for each $D_i$, of a permutation $\sigma$ of the grades preserving their ranks and digits $c_j\in\{0,\ldots,p-1\}$ nondecreasing in the source order. We also require $\sigma$ to be increasing on every block on which the digit $c_j$ is constant. In local adapted bases \citep[Proposition~3, pp.~11--12]{Borne2007D} $v_j$ and $w_k$ and a local equation $t$ of $D_i$, we impose on the matrix $(a_{kj})$ of a map $F^*E\to E'$ the conditions
\[v_t(a_{kj})\ge c_j+\mathbf1_{k<\sigma(j)},\]
and require the diagonal blocks $k=\sigma(j)$, after division of column $j$ by $t^{c_j}$, to be invertible modulo $t$. Away $D$ the map is required to be an isomorphism.

These local conditions are intrinsic. The normalized Frobenius source lattice is
\[L_c=\sum_j\mathcal O(c_jD_i)\,F^*E_{i,j-1}.\]
Locally it has basis $t^{-c_j}F^*v_j$. For an adapted source change of basis, an entry from a higher-weight column to a lower-weight row is divisible by $t$; after Frobenius it is divisible by $t^p$. Since digit differences are at most $p-1$, the normalized transition matrix is integral. Modulo $t$, entries between unequal digit blocks vanish, while within each digit block the original flag order is preserved. The digit blocks and their subflags therefore define a canonical reordered flag on $L_c|_{D_i}$. The displayed condition is precisely an isomorphism of this flagged bundle with the target flagged bundle.
All these constructions use only finitely many sums, images, restrictions to fixed thickenings, and twists of order at most $p-1$. Thus they are constructible in the bounded families. No root stack is extended across the bad points: maps of reflexive lattices extend uniquely across $\Sigma$, and on $Y$ their poles are controlled by the fixed exceptional-window bound. The finitely many digit twists add only a fixed amount to that bound. We obtain finitely many constructible correspondences
\[R_c\subset P_r\times P_r.\]
Consider a closed word of $f$ such types. Following a grade around the induced permutation cycles, solve
\[\alpha_{\mathrm{next}}=p\alpha-c.\]
On a cycle of length $h$ there is a unique solution, with denominator dividing $p^h-1$. Here $h\le fr$. We discard words for which the solutions fail to lie in $[0,1)$, fail to have the prescribed strict order, or identify two grades prescribed to be distinct. In particular, the all-$(p-1)$ cycle, which gives weight $1$, is excluded. There are at most $A_r^f$ words before this discard.
For a valid word choose a common denominator $n$ prime to $p$. The flagged states become parabolic vector bundles on $(S^\circ,D)$ with these weights, hence vector bundles on the root stack $\mathcal X=\sqrt[n]{(S^\circ,D)}$ by Borne's equivalence \citep[Theorem~4, p.~12, with proof through p.~16]{Borne2007D}. This equivalence is being used only where the boundary components are smooth and disjoint. The edge criterion is exact: if the source and target weights are $\alpha_j$ and $\beta_k$, then
\[\left\lceil p\alpha_j-\beta_k\right\rceil=c_j+\mathbf1_{k<\sigma(j)}.\]
This is exactly the condition for the pulled-back matrix to be regular on the Kummer chart. Its normalized diagonal blocks are invertible precisely when its determinant is a unit there. Hence the map and its inverse are both regular on that chart. Uniqueness on the dense complement glues these maps to genuine Frobenius isomorphisms on $\mathcal X$; the parabolic gluing and tensor compatibility are also explicit in \citet[Theorem~6.1, pp.~33--36]{BorneVistoli2012D}.
Composing the $f$ edges gives a $p^f$-unit vector bundle on $\mathcal X$. Its fixed-frame equations, locally $z^{[p^f]}=C^{-1}z$, define a finite étale $\mathrm{GL}_r(\mathbf F_{p^f})$-torsor. Its restriction to $U^\circ$ has inertia factoring through $\mu_n$ at every original boundary component, so it defines a representation of $Q^\circ$. Conversely, every such representation gives a valid word and a cycle of states.
For a fixed weighted state, its algebra of global parabolic endomorphisms is a finite-dimensional subalgebra of $H^0(S,\mathcal End(E))$. This follows from reflexive extension across the finite omitted set and from preservation of the finitely many lattices. A given $p^f$-unit structure induces a $p^f$-semilinear endomorphism $\tau$ of this algebra. Its unit group is connected. The $\tau$-twisted conjugation action has open orbits, because $d\tau=0$ and the differential in the group variable is invertible. It therefore has one orbit. Thus a weighted state admits at most one unit structure up to isomorphism; choices of intermediate edge isomorphisms do not produce extra representation classes.

Conversely, for a fixed representation and fixed valid word, its finite étale extension to the corresponding root stack is unique: restriction from a normal root stack to its dense open complement is fully faithful, as is seen on its normal Kummer charts. The associated unit bundle and its entire parabolic chain are consequently unique, and all intermediate states are its Frobenius conjugates. Thus the cycle set for a fixed valid word is in bijection with a subset of the representation classes just described. The fixed-original-boundary smallness theorem proved earlier, Theorem~\ref{D:fixed-boundary-small}, makes this set finite. This assertion is made only for valid words.
Choose finite basic affine locally closed descriptions of the fixed constructible correspondences $R_c$. For a word of length $f$, choosing one such piece at each edge gives a cycle locus in $O_r(f)$ variables, defined by $O_r(f)$ equations and inequations of degrees bounded independently of $f$. Introduce an inverse variable for each inequation. Its geometric point set is finite, so the resulting affine scheme is zero-dimensional after reduction. The elementary Bézout bound gives at most $\exp(C_r f)$ points. The number of piece words is also at most exponential in $f$. Summing proves the asserted bound for $Q^\circ$, and hence for its quotient corresponding to the original omitted set.
We emphasize that this argument uses a finite collection of correspondences. It does not require a single-valued Frobenius function on the unweighted state quotient. It also does not provide a rank-independent bound for the constants $C_r$, and therefore does not settle the remaining nonabelian finite-generation problem.
\end{proof}
